%% notes-tex/database/kg-build-system/sections/3-release-pairing.tex

\section{Releases and pairing\secstatus{todo}}
\label{sec:pairing}

%% SOURCE: notes/kg-build-system.mermaid.release-lifecycle.md, rendered by
%% notes/assets/publish/scripts/mermaid_pdf.sh and scaled to the text block by
%% `make diagrams`; the gates are those of PR #680 (notes/versioning.md, 2026.10.06).
\tcfig{figures/release-lifecycle.pdf}{The life of one release. A package tag is
cut first, the build records it, the snapshot commit bumps the package again, and
from then on every reader of the store is checked against the contracts the store
was written with. The red hexagons are the four gates; each one refuses rather
than degrades.}{fig:lifecycle}

\notesec{Why a pair}
Every served record is serialized under the schema contract of the commit that
built it, and the snapshot records, per dataset, the fingerprints of that
dataset's schema closure. A package version reads a release correctly only when
its schema surface reproduces every one of those fingerprints. The October 2026
builds showed what happens without a rule: both ran from an untagged \texttt{main},
their snapshot commits used a tag the release parser does not bump on, and the
compatibility page then showed every published package as partially compatible
with the served store while the client read it anyway. A pair is the rule that
replaces "partially": a release names one package, and a package either reads a
release in full or is refused.

\notesec{The four gates}
\begin{enumerate}
\item \textbf{Build.} The slurm script refuses to start unless the build commit
  carries a \texttt{vX.Y.Z} tag and the fingerprint checkout sits on it, so the
  stamp records the tag. A release is therefore cut before a build: a \texttt{REL}
  (minor) or \texttt{DB}/\texttt{PATCH} (patch) commit on \texttt{main},
  semantic-release's bump commit and tag, then the build from that commit.
\item \textbf{Snapshot.} The build writes \texttt{database/releases/<release>.json}
  and its closures file into the checkout; they are committed as
  \texttt{DB(kg): ...}, a patch bump, so the package line advances with the store.
\item \textbf{Client.} At connect, \texttt{Neo4jQueryRaw} reads the resolved
  database's \texttt{KgRelease} node and runs the pairing check
  (\texttt{require\_paired} in \file{torchcell/knowledge_graphs/releases.py})
  against the installed schema surface before opening a driver. Any drifted or
  unverified dataset, or a store with no release node, raises
  \texttt{IncompatibleReleaseError} naming the paired package and the install
  command. There is no override.
\item \textbf{Page.} \file{scripts/kg_compat_page.py} renders a pairs table and
  the tag-by-release matrix from the committed snapshots; a paired tag that stops
  reading its release is a broken pair and the generator refuses, which fails the
  \texttt{--check} step in CI.
\end{enumerate}

%% SOURCE: hand-authored from docs/source/database/compatibility.md as generated on
%% 2026-10-06 by scripts/kg_compat_page.py.
\begin{table}[htbp]
\centering
\footnotesize
\caption{The pairs as of 2026-10-06. "Reads it" lists every published tag whose
schema surface reproduces all 51 closures; the paired package is the one the
release names.}
\label{tab:pairs}
\begin{tabular}{llll}
\toprule
KG release & KG version & paired package & reads it \\
\midrule
\texttt{2026.09.21-ab6d8c5d} & 1.2 & v1.2.1 & v1.2.1 to v1.6.1 \\
\texttt{2026.10.02-833970cd} & 2.0 & v1.6.1 & v1.2.1 to v1.6.1 \\
\texttt{2026.10.06-4b293d34} & 3.0 & v1.6.2 & v1.6.2 \\
\bottomrule
\end{tabular}
\end{table}

\notesec{Repairing a release built before the rule}
\texttt{releases retag --release <id> --tag vX.Y.Z} pairs a snapshot built from an
untagged commit with a tag cut afterwards, and accepts the tag only when the
schema surface at that tag reproduces every served closure. It is idempotent for
the same tag and refuses a second, different tag. \texttt{kg\_release.sh retag}
runs it with the manifest and then rewrites the served node. That is how the
three releases in Table~\ref{tab:pairs} got their packages; from the next build
on, the stamp records the tag and no repair is needed.
